\section{Introduction}

This report documents both the boot-process background behind operating-system startup and the development of a small 32-bit operating-system kernel for the \texttt{OSDev\_18} project. It starts by explaining the firmware and bootloader stages that happen before a kernel runs, then follows the project implementation from a minimal bootable kernel into a more interactive environment with processor setup, interrupts, memory management, timing, sound, keyboard input, and simple applications.

The background section covers the general boot sequence: Power-On Self-Test, firmware handoff, BIOS and UEFI differences, bootloader responsibilities, early i386 memory layout, modern operating-system boot chains, and the role of virtualization. This gives context for why the project can rely on a bootloader, why the kernel starts from a constrained early machine state, and why QEMU is useful for testing this kind of low-level work.

The implementation itself is intentionally low-level. Most of the kernel interacts directly with x86 processor structures and hardware interfaces, including the Global Descriptor Table, Interrupt Descriptor Table, Programmable Interrupt Controller, Programmable Interval Timer, VGA text output, keyboard controller, and PC speaker. Because the kernel runs in a freestanding environment, basic facilities that are normally provided by an operating system or standard library also had to be implemented inside the project.

After the boot background, the report follows the same order as the implementation work. The first implementation part covers the early boot path and processor setup, with particular focus on the Global Descriptor Table and the need to reload segment registers after installing it. The second part describes interrupt handling, including CPU exceptions, hardware IRQs, PIC remapping, and keyboard input. The third part introduces memory management, paging, PIT-based timing, sleep functions, and PC speaker music playback. The final part shows how these individual subsystems were combined into a small menu-driven application environment with a music player and an interactive Snake game.

The overall goal was not to build a complete general-purpose operating system, but to demonstrate how core operating-system mechanisms fit together. Each stage adds a specific capability, and later stages reuse earlier ones: the PIT depends on interrupt handling, the music player depends on PIT timing and PC speaker output, and Snake depends on memory allocation, keyboard input, terminal drawing, timing, and sound feedback. By the end of the project, the kernel has moved from static startup output to an event-driven system that can run simple interactive applications.
